\addtocontents{toc}{\protect\Needspace{3\baselineskip}}
\originalchapter{8.821 课程详讲：D3 喉部与 AdS 几何}
本章对应第 15 讲末、第 16 讲和第 17 讲的几何部分. 先使用已知 Type IIB D3 解, 复核其尺度、调和函数和近视界极限, 再解释低能解耦为何仍留下完整的喉部弦态. Einstein 张量的一般计算以前置 GR 卷为依据; 不把验证一个调和函数当作已经重推全部 IIB 超引力方程.

\originalsection{一张带 RR 荷的膜怎样产生几何}
点电荷在三维空间有 $1/r$ 静电势, 相应 Einstein--Maxwell 源能产生 Reissner--Nordström 几何. D3 膜在十维中有六个横向空间方向, 同样同时携荷和张力, 它的引力与五形式场强共同改变背景. 与点电荷比较的作用是说明源、规范场能量和几何须共同满足场方程, 不是把十维膜强行塞进四维点粒子解.

$N$ 张重合、平直 BPS D3 的标准解是
\[
\dd s^2=H(r)^{-1/2}\eta_{\mu\nu}\dd x^\mu\dd x^\nu
+H(r)^{1/2}(\dd r^2+r^2\dd\Omega_5^2),
\quad H(r)=1+\frac{R^4}{r^4},
\quad R^4=4\pi g_sN(\alpha')^2.
\]
$\mu=0,1,2,3$, $r$ 是六个法向坐标的半径. $R$ 是由源决定的长度尺度, 不是输入的宇宙半径. 离开源处球对称的六维调和方程为
\[
\nabla_6^2H=\frac1{r^5}\frac\dd{\dd r}
\left(r^5\frac{\dd H}{\dd r}\right)=0.
\]
积分给常数加 $r^{-4}$, 由渐近平直取常数为1. 这个检查说明幂次为何是4; BPS 关系、五形式与各个度规幂次仍来自 IIB 的实际引力解.

在一种常用 string-frame RR 归一化中, 五形式可写
\[
F_5=\frac1{g_s}(1+*)\left[\dd(H^{-1})\wedge
\dd t\wedge\dd x^1\wedge\dd x^2\wedge\dd x^3\right].
\]
适当取方向后, 球面磁通为 $\int_{S^5}F_5=4\pi^3R^4/g_s$. 量子化 $\int F_5=(2\pi)^4(\alpha')^2N$ 即给 $R^4=4\pi g_sN(\alpha')^2$. 这个版本的 $F_5$ 与前章动能标准归一化荷的场变量有常数重标度, 必须沿一套约定计算, 而不能只比较一个叫作“电荷”的数字.

当 $r\gg R$, $H\simeq1$, 几何近十维平直, 微小改变随 $R^4/r^4$ 衰减. 当 $r\sim R$, 几何变形不可忽略, 曲率量级 $R^{-2}$. 为让导数修正小, 要求 $\alpha'/R^2\ll1$, 即 $g_sN\gg1$. 还要压低微扰弦环修正, 通常取 $g_s\ll1$. 因此有用的超引力背景不是只令耦合弱而固定少数膜, 而是弱 $g_s$ 与大量 $N$ 同时实现.

\originalsection{近视界乘积由每个系数直接读出}
在 $r\ll R$, 可忽略 $H$ 的1, 得
\[
\dd s^2=\frac{r^2}{R^2}(-\dd t^2+\dd\mathbf x^2)
+\frac{R^2}{r^2}\dd r^2+R^2\dd\Omega_5^2.
\]
最后一项的球半径\textbf{不再随 $r$ 改变}, 是固定 $R$ 的 $S^5$. 前两项令 $z=R^2/r$, 则变成 $R^2(-\dd t^2+\dd\mathbf x^2+\dd z^2)/z^2$, 是 $\mathrm{AdS}_5$. 因此喉部几何是 $\mathrm{AdS}_5\times S^5$, 两因子曲率半径相同.

沿定时间径向片, 喉部固有距离为 $\int R\dd r/r$, 向 $r=0$ 对数发散. 源不再表现为某个可到达的普通平直空间点; 它的信息由喉部几何和球面通量承载. 这不能由无限空间距离推出自由落体永不能进入, 因为时间型轨迹的固有时与定时间空间长度不同.

\originalsection{低能极限为何在两种描述中留下同样部门}
开放弦描述中, 重开弦模式质量约 $1/\sqrt{\alpha'}$, 低能 $E\sqrt{\alpha'}\to0$ 留下膜上的无质量场. D3 的有效理论是四维 $\mathcal N=4$ $U(N)$ SYM. 渐近平直的无质量闭弦则成为自由引力等场, 因为 $G_{10}\propto g_s^2(\alpha')^4$ 使其与给定能标的开放部门相互作用消失. 这一步的低能截断在固定边界能量下进行.

引力描述必须使用同一个远处时间 $t$ 来比较能量. 固定 $r$ 的局部固有时为 $\dd\tau=H^{-1/4}\dd t$, 故
\[
E_\infty=H^{-1/4}E_{\rm local},\qquad
E_{\rm local}=H^{1/4}E_\infty.
\]
近喉部 $H^{-1/4}\simeq r/R$, 一份局部很高能的弦态也可能在远处显示有限低能. 局部能量乘错红移因子, 就会错误地把喉部全部重弦也删掉.

可用实际尺度极限看清这一点. 固定 $U=r/\alpha'$、$g_s,N$ 和 $\lambda=4\pi g_sN$, 令 $\alpha'\to0$. 近喉部有
\[
\frac{\dd s^2}{\alpha'}=
\frac{U^2}{\sqrt\lambda}(-\dd t^2+\dd\mathbf x^2)
+\sqrt\lambda\left(\frac{\dd U^2}{U^2}+\dd\Omega_5^2\right).
\]
几何在弦长度单位下保持非平凡. 若局部弦态 $E_{\rm local}\sim1/\sqrt{\alpha'}$, 则
\[
E_\infty\sim\frac rR\frac1{\sqrt{\alpha'}}
=\frac{U}{\lambda^{1/4}}
\]
仍有限. 所以喉部保留\textbf{完整弦理论}, 只有在进一步 $\lambda\gg1$ 时才可把其中重弦在低曲率过程里近似积分掉. 固定 $r$ 的平直远区重弦则确实退出低能谱.

两种描述各自都留下一个相同的平直自由闭弦部门. 去掉这一共同部门, 把剩下的膜 SYM 与喉部 IIB 弦理论联系起来, 就得到 Maldacena 对偶猜想的构造动机. 解耦结构、对称和后续检验为它提供依据, 并没有从上述有限几行尺度变换完成任意 $N,\lambda$ 的普适证明.

还有一个常见疑问: 膜上的有效 Lagrangian 只含无质量基本场, 为何能对偶于有重弦态的体理论? 场论的局域复合算符谱不限于 Lagrangian 中的基本场; 在强耦合下它们可有很大的标度维数. 在全局边界球面上, 标度维数对应能量, 体的重弦对应高维数复合态. 后面的质量--维数关系把这个联系定量化.

\originalsection{AdS 曲率和 Poincaré patch}
一般 $\mathrm{AdS}_{d+1}$ 取
\[
\dd s^2=\frac{R^2}{z^2}
(\dd z^2+\eta_{\mu\nu}\dd x^\mu\dd x^\nu),\qquad z>0.
\]
其非零 Christoffel 结构可由 $z$ 的单一依赖快速复核:
\[
\Gamma^z{}_{zz}=-\frac1z,\quad
\Gamma^z{}_{\mu\nu}=\frac{\eta_{\mu\nu}}z,\quad
\Gamma^\mu{}_{z\nu}=-\frac{\delta^\mu{}_{\nu}}z.
\]
代入前置 GR 的曲率定义得
\[
R_{MN}=-\frac d{R^2}g_{MN},\quad
\mathcal R=-\frac{d(d+1)}{R^2},\quad
\Lambda=-\frac{d(d-1)}{2R^2}.
\]
由此得 $R_{MNPQ}=-(g_{MP}g_{NQ}-g_{MQ}g_{NP})/R^2$, 是恒负曲率. 本节用 $\mathcal R$ 写 Ricci 标量, 避免与半径 $R$ 同符号引起量纲错读.

AdS 可先表示为签名 $(2,d)$ 的嵌入双曲面 $-X_{-1}^2-X_0^2+\sum_{i=1}^dX_i^2=-R^2$. Poincaré 映射的一套明确公式为
\[
X_{-1}=\frac{R^2+z^2+x^2}{2z},\quad
X_d=\frac{R^2-z^2-x^2}{2z},\quad
X_\mu=\frac{Rx^\mu}{z},
\]
其中 $x^2=\eta_{\mu\nu}x^\mu x^\nu$, $X_0=Rt/z$. 代回可检查双曲面方程; $X_{-1}+X_d=R^2/z>0$, 所以它只覆盖这一半空间条件对应的 patch. $z=\infty$ 是 patch 视界, 并不是 AdS 的曲率奇点.

\originalsection{全局时间和共形边界}
用无量纲全局半径 $s\ge0$、单位球方向 $n_i$, 写
\[
X_{-1}=R\sqrt{1+s^2}\cos t_G,\quad
X_0=R\sqrt{1+s^2}\sin t_G,\quad
X_i=Rsn_i.
\]
诱导度规为
\[
\dd s_{
\rm geom}^2=R^2\left[-(1+s^2)\dd t_G^2+
\frac{\dd s^2}{1+s^2}+s^2\dd\Omega_{d-1}^2\right].
\]
嵌入双曲面原本让 $t_G$ 周期为 $2\pi$, 会有闭合时向曲线. 物理应用通常取其普适覆盖, 将 $t_G$ 展开为实数. 上述嵌入表达是每个时间分支的局部表示, 不再用重复的嵌入坐标把不同覆盖点误认成同一点.

令 $s=\tan\varrho$, $0\le\varrho<\pi/2$, 得
\[
\dd s_{\rm geom}^2=\frac{R^2}{\cos^2\varrho}
[-\dd t_G^2+\dd\varrho^2+\sin^2\varrho\dd\Omega_{d-1}^2].
\]
乘去正共形因子, 边界 $\varrho=\pi/2$ 的度规是 $-\dd t_G^2+\dd\Omega_{d-1}^2$, 即 $\mathbb R\times S^{d-1}$. 半径为 $R_b$ 的边界代表取 $t_b=R_bt_G$, 不能将这个共形代表半径无说明地与体 $R$ 当成同一物理尺.

径向零曲线满足 $\dd t_G=\pm\dd\varrho$, 从中心到边界所需坐标时为 $\pi/2$. 有质量自由落体有限能量 $e$ 则由归一化得到 $(\dd s/\dd(\tau/R))^2=e^2-1-s^2$, 最远 $s=\sqrt{e^2-1}$, 不到边界. 因而全局 AdS 对这类激发像一个有限大小的引力盒, 合适边界条件下能谱离散, 而不是靠在某处放硬墙模拟的平直盒.

\begin{example}
取 $e=\sqrt2$, 径向粒子最大到 $s=1$. 解 $s=\sin(\tau/R)$, 从中心出发到再次回中心经历固有时 $\pi R$. 对全局时间, 积分 $\dd t_G/\dd(\tau/R)=e/(1+s^2)$ 也给返回间隔 $\pi$. 这个与振幅无关的返回时间在 PS4.2 中完整考核; 它不表示两个具有不同绕行数的覆盖空间事件都可仅由嵌入弦长判断因果关系.
\end{example}

\originalsection{等距群为何匹配共形群}
嵌入度规保持群是 $SO(2,d)$. 在 Poincaré 坐标中, 它包含 $d$ 个平移、$d(d-1)/2$ 个 Lorentz 变换、一个联合 $x,z$ 尺度变换和 $d$ 个特殊共形变换. 总数 $(d+1)(d+2)/2$, 与边界共形群一致.

特殊共形的一套形式是
\[
z'=\frac zD,\qquad x'^\mu=\frac{x^\mu+b^\mu(z^2+x^2)}D,
\quad D=1+2b\cdot x+b^2(z^2+x^2).
\]
它可由反演、平移、再次反演构造. 在想保持同一个 $z>0$ patch 时需 $D>0$; 零分母或变号表示 patch 条件失效, 不能称为全局没有任何定义域问题的普通分式变换. PS4.1 给出度规不变的具体计算. 对称性匹配是对偶必要的强检验, 但单独一个相同的群仍不能固定所有相关函数的动力学常数.
